\pageoriginale

\clearpage
\SGAExposeBookmark{XV}{La fórmula de Picard--Lefschetz}
\begin{center}
{\large SGA 7 --- EXPOSÉ XV}\par
\vspace{2em}
{\Large\underline{La fórmula de Picard--Lefschetz}}\par
\vspace{1.5em}
{\large por P. Deligne}
\end{center}

\medskip
\noindent\textbf{Contenido de la fuente}
\begin{center}
\begin{tabular}{@{}llr@{}}
\S1. & Singularidades cuadráticas ordinarias (formas canónicas) & 2\\
1.1 & Recordatorios & 2\\
1.2 & Singularidades cuadráticas ordinarias & 4\\
1.3 & Módulos de singularidades cuadráticas & 10\\
\S2. & Cálculo de los ciclos evanescentes en un caso cuadrático
  ordinario estándar & 13\\
2.1 & Preliminar: cohomología de un cono & 13\\
2.2 & Ciclos evanescentes & 17\\
\S3. & La fórmula de Picard--Lefschetz & 23\\
3.1 & Resultados preliminares & 23\\
3.2 & Dimensión relativa par & 24\\
3.3 & Dimensión relativa impar & 26\\
3.4 & Resumen & 31
\end{tabular}
\end{center}

\pageoriginale

\section*{1. Singularidades cuadráticas ordinarias (formas canónicas)}

En esta sección retomamos los resultados de VI 6, adoptando un punto de
vista ``henseliano'' en lugar de ``formal''. Damos formas canónicas para
las singularidades cuadráticas ordinarias y sus deformaciones, en las
que pueden ponerse localmente para la topología étale.

Estos resultados se deducen fácilmente de VI 6, del teorema de la
función implícita (en la forma de Tougeron) y del teorema de R. Elkik
sobre los módulos henselianos versales de singularidades aisladas. En
el momento del seminario oral, este último aún no estaba disponible;
se había demostrado un teorema más débil, que cubría únicamente el caso
de las intersecciones completas (lo cual nos bastaba), y que no se
publicará.

\subsection*{1.1. Recordatorios}

\subsubsection*{1.1.1.}
Sea $f:X\longrightarrow S$ un morfismo, con $X$ localmente descriptible
por $p$ ecuaciones en el espacio afín relativo $\mathbb A_S^{n+p}$ de
dimensión $n+p$ sobre $S$. El ideal jacobiano $J^n(X/S)$ se definió en
VI 5. Si $X$ es el subesquema cerrado de $\mathbb A_S^{n+p}$ definido
por las ecuaciones $f_i=0$ $(1\leq i\leq p)$, es el haz de ideales
engendrado por los menores de orden $p\times p$ de la matriz jacobiana
$(\partial f_i/\partial x_j)$. El teorema siguiente es bien conocido.
Un caso particular de su análogo formal se demuestra en Bourbaki,
\emph{Alg\`ebre commutative}, III, \S4, n\textsuperscript{o}~6,
Corolario 3 del Teorema 2; el análogo $C^\infty$ se demuestra en
Tougeron, \emph{Id\'eaux de fonctions diff\'erentiables} (tesis,
Rennes, mayo de 1967); el caso considerado aquí se trata en Artin,
\emph{Algebraic approximation of structures over complete local
rings}, Publ. Math. IHES \underline{36} (volumen de Zariski), Lema 5.10.

\subsubsection*{Teorema 1.1.2 (Tougeron--Artin).}
\underline{Sea $S$ el espectro de un anillo local henseliano $A$, sea
$f:X\to S$ un morfismo como en 1.1.1, sean $a$ e $I$ dos ideales
propios de $A$, y sea $\delta'$ un haz de ideales de $X$ contenido en
$J^n(X/S)$. Supongamos que $I$ es de tipo finito. Sea $s$ una sección
de $f$ sobre $S/I$, es decir, $s\in X(A/I)$. Denotemos también por
$s^*\delta'^2$ el ideal de $A$ imagen inversa del ideal
$s^*\delta'^2$ de $A/I$.}

\underline{Si $(s^*\delta'^2)\,a\supset I$, existe una sección
$s_1\in X(A)$ de $f$}

\pageoriginale

\underline{congruente con $s$ módulo $(s^*\delta')\,a$.}

\subsubsection*{1.1.3.}
Sea $S$ el espectro de un anillo local henseliano, de punto cerrado
$s$. Sea $X'_0$ un esquema de tipo finito sobre $k(s)$, liso fuera de
un punto cerrado $x_0$, y sea $X_0$ el henselizado de $X'_0$ en $x_0$.
Nos interesan los sistemas
\begin{center}
\begin{tikzcd}[column sep=large,row sep=large]
X_0 \arrow[r,hook] \arrow[d]
  & X \arrow[d]
  & {} \\
s \arrow[r,hook]
  & T \arrow[r]
  & S
\end{tikzcd}
\end{center}
del tipo siguiente:
\begin{enumerate}[label=(\alph*),leftmargin=2.6em]
\item $T$ es local henseliano y esencialmente de presentación finita
sobre $S$ (I 0.0.5); el morfismo $T\to S$ es local e induce un
isomorfismo de cuerpos residuales;

\item $X$ es local henseliano, plano y esencialmente de presentación
finita sobre $T$;

\item la fibra especial de $X/T$ está dotada de un $k(s)$-isomorfismo
con $X_0$.
\end{enumerate}
Estas son las \underline{$S$-deformaciones henselianas} de $X_0$. Un
\underline{morfismo} de $(X,T)$ a $(X',T')$ es un diagrama conmutativo
de $S$-esquemas
\begin{center}
\begin{tikzcd}[column sep=large,row sep=large]
X \arrow[r] \arrow[d] & X' \arrow[d] \\
T \arrow[r] & T'
\end{tikzcd}
\end{center}
que induce la identidad de $X_0$ al pasar a las fibras especiales. Una
deformación $X/T$ se llama \underline{versal} (minimalmente versal) si:
\begin{enumerate}[label=(\alph*),leftmargin=2.6em]
\item para todo morfismo $X''/T''\longrightarrow X'/T'$ para el cual
$T''\to T'$ es una inmersión cerrada, la aplicación
\[
\operatorname{Hom}(X'/T',X/T)
\longrightarrow\operatorname{Hom}(X''/T'',X/T)
\]
es sobreyectiva. En particular, tomando $X''/T''=X_0/s$, se tiene
$\operatorname{Hom}(X'/T',X/T)\ne\varnothing$;

\item para $T'=\operatorname{Spec}(k(s)[\varepsilon]/(\varepsilon^2))$,
la aplicación de ``imagen inversa''
\[
\operatorname{Hom}(T',T)\longrightarrow
\{\text{clases de isomorfismo de deformaciones de $X_0$ sobre $T'$}\}
\]
es biyectiva.
\end{enumerate}

\pageoriginale

\subsubsection*{Teorema 1.1.4 (R. Elkik).}
\underline{Existe una $S$-deformación henseliana versal de $X_0$ (y es
única salvo isomorfismo no único). Si $S$ es noetheriano, una
deformación $X/T$ es de este tipo si y solo si $X$ induce sobre la
completación $\widehat T$ de $T$ una $S$-deformación formal versal de
$X_0$ (una deformación formal minimalmente versal; VI 1.9 y VI 4).}

Para la demostración de este delicado teorema, remitimos a R. Elkik
(seminario de la ENS, 1971/72).

\subsection*{1.2. Singularidades cuadráticas ordinarias}

Sea $y$ un punto cerrado de un esquema $Y$ de tipo finito sobre un
cuerpo $k$ de característica $p$, y sea $n$ la dimensión de $Y$ en $y$.

\subsubsection*{Definición 1.2.1 (= VI 6.6).}
\begin{enumerate}[label=(\roman*),leftmargin=2.6em]
\item \underline{Si $k$ es algebraicamente cerrado, se dice que $y$ es
un punto cuadrático ordinario de $Y$ si la completación
$\widehat{\mathcal O}_{Y,y}$ del anillo local de $Y$ en $y$ es isomorfa
al cociente de $k[[x_1,\ldots,x_{n+1}]]$ por el ideal engendrado por una
única serie formal $f$ que puede escribirse}
\begin{equation}\tag{\sisi{4.2.1}{1.2.1.1}}
f(\underline X)=Q(\underline X)\sisi{=}{+}
\text{términos de orden }>2,
\end{equation}
\underline{donde $Q(\underline X)$ es una forma cuadrática ordinaria.}

\item \underline{Si $\bar k$ es una clausura algebraica de $k$, se dice
que $y$ es un punto cuadrático ordinario de $Y$ si los puntos de
$\bar Y=Y\otimes_k\bar k$ situados encima de $y$ son puntos cuadráticos
ordinarios de $\bar Y$.}
\end{enumerate}

\subsubsection*{1.2.2.}
Reemplazando ``forma cuadrática ordinaria'' por ``forma cuadrática no
degenerada'' en esta definición, se obtiene la noción de
\underline{punto cuadrático no degenerado} (VI 6.0). Un punto $y$ es un
punto cuadrático no degenerado de $Y$ si y solo si es un punto
cuadrático ordinario y o bien $p\ne2$, o bien

\pageoriginale

$n$ es impar.

\subsubsection*{Ejemplo 1.2.3.}
Sea $Q(\underline X)$ una forma cuadrática no degenerada en $n+1$
variables sobre $k$. Sea $Y_0$ el cono cuadrático de
$\mathbb A_k^{n+1}$ de ecuación $Q=0$. Entonces el origen, denotado por
$y_0$, es un punto cuadrático no degenerado de $Y_0$.

\subsubsection*{Ejemplo 1.2.4.}
Supongamos que $p=2$ y $n=2m$. Sea $a\in k$, y sea $Q$ el polinomio
\[
Q(\underline X)=(X_0^2-a)
 +\sum_{0<i\leq j\leq2m}a_{ij}x_ix_j.
\]
Supongamos que la forma cuadrática en $2m$ variables
$\sum_{0<i\leq j}a_{ij}x_ix_j$ es no degenerada, y denotemos por $y_0$
el punto cerrado de $\mathbb A_k^{n+1}$ de coordenadas
$(\sqrt a,0,\ldots,0)$. Entonces $y_0$ es un punto cuadrático ordinario
del subesquema cerrado $Y_0\subset\mathbb A_k^{n+1}$ de ecuación $Q=0$.

\subsubsection*{Ejemplo 1.2.5.}
Sea $k'$ una extensión separable finita de $k$, y sea $Y_0$ un
$k'$-esquema del tipo considerado en 1.2.3 (respectivamente, 1.2.4). Si
$Y_0$ se considera como $k$-esquema, entonces $y_0$ es un punto
cuadrático no degenerado (respectivamente, ordinario) de $Y_0$. La
extensión $k'/k$ se identifica con la mayor subextensión separable de
$k(y_0)/k$.

\subsubsection*{Teorema 1.2.6.}
\underline{Supongamos que $y$ es un punto cuadrático ordinario del
$k$-esquema $Y$, y sea $k'$ la mayor subextensión separable de
$k(y)/k$. Existen entonces un $k'$-esquema
$Y_0\subset\mathbb A_{k'}^{n+1}$ del tipo de 1.2.3 o 1.2.4 y un
$k$-isomorfismo $\varphi$ entre el henselizado $Y_{(y)}$ de $Y$ en $y$
y el henselizado $Y_{0(y_0)}$ de $Y_0$ en $y_0$.}

Obsérvese primero que existe un entorno étale $Y'$ de $y$ en $Y$
\begin{center}
\begin{tikzcd}[row sep=large,column sep=large]
{} & Y'\arrow[d]\\
y\arrow[ur,hook]\arrow[r,hook]&Y
\end{tikzcd}
\end{center}
que es un $k'$-esquema. En efecto, $Y_{(y)}$ es canónicamente un
$k'$-esquema.

\pageoriginale

Reemplazando $(k,Y,y)$ por $(k',Y',y)$, nos reducimos al caso
$k=k'$, es decir, al caso en que $k(y)/k$ es una extensión radicial.

El $k(y)$-espacio vectorial $(\Omega^1_{Y/k})_y$ tiene dimensión
$n+1$. Basta comprobarlo después de extender los escalares de $k(y)$
a su clausura algebraica y completar en $y$, donde es inmediato. Por
tanto, un entorno de $y$ en $Y$ es un subesquema cerrado de un
$k$-esquema liso $Z$ de dimensión $n+1$. Además, reemplazando $Y$ y
$Z$ por entornos adecuados de $y$, el ideal de $\mathcal O_Z$ que
define $Y$ es principal (por el mismo argumento anterior), digamos de
ecuación $f$.

\subsubsection*{Lema 1.2.7.}
\underline{Supongamos que $y$ es un punto cuadrático no degenerado de
$Y$ y que $k(y)/k$ es radicial. Entonces, en un entorno de $y$,
$J^n_{Y/k}$ (1.1.1) está engendrado por un elemento, soportado en $y$,
y es de rango uno sobre $k$. En particular, $k(y)=k$.}

Basta demostrar la primera afirmación después de pasar a una clausura
algebraica de $k$ y completar en $y$. Entonces es inmediata por cálculo
directo: con la notación de 1.2.1, el lema de Nakayama muestra que los
$\partial f/\partial x_i$ (y por tanto, con mayor razón, $f$ y los
$\partial f/\partial x_i$) engendran el ideal maximal.

Demostremos 1.2.6 cuando $y$ es no degenerado. Podemos y vamos a
suponer que $k(y)=k$, que $Y$ está definido por una ecuación $f$ en un
esquema liso $Z$, y que existe un morfismo étale
\[
x:Z\longrightarrow\mathbb A_k^{n+1}
\]
(un ``sistema de coordenadas locales $x_i$'') que lleva $y$ a $0$.
Existe entonces una forma cuadrática no degenerada $Q$ en $n+1$
variables tal que, en $y$,
\[
f=Q(x_i)+\text{términos de orden }>2.
\]
Sea $Y_0$ el esquema definido por $Q$ como en
\sisi{4.3}{1.2.3}.

\underline{Sobre $Y$}, el sistema $(x_i)$ es una ``solución aproximada''
de la ecuación $Q(x_i)=0$: si $\mathfrak m$ es el ideal que define
$y$, entonces
\[
Q(x_i)\in\mathfrak m^3\qquad\text{sobre }Y.
\]

\pageoriginale

Como $Q$ es no degenerada, el ideal de $\mathbb A_k^{n+1}$ engendrado
por los $\partial Q/\partial x_i$ es el ideal maximal en $0$, y su
imagen inversa por $x=(x_i):Y\to\mathbb A_k^{n+1}$ es $\mathfrak m$.
Por el teorema de la función implícita 1.1.2 (aplicado con
$a=\mathfrak m$ y $\delta'=J$), existen sobre el henselizado $Y_{(y)}$
funciones $x'_i$, congruentes con los $x_i$ módulo $\mathfrak m^2$,
que satisfacen $Q(x'_i)=0$. Proporcionan el isomorfismo buscado entre
$Y_{(y)}$ y $Y_{0(y_0)}$.

\subsubsection*{Lema 1.2.8.}
\underline{Supongamos que $y$ es un punto cuadrático ordinario
degenerado ($p=2$ y $n$ par), y que $k(y)/k$ es radicial. Entonces, en
un entorno de $y$, $T^1_{Y/k}$ está engendrado por un elemento,
soportado en $y$, y es de rango dos sobre $k$. En particular, o bien
$k(y)=k$, o bien $k(y)$ es isomorfo a $k(\sqrt a)$ para algún
$a\in k$, $a\notin k^2$.}

Procedamos como en 1.2.7. Con la notación de 1.2.1, basta comprobar que
\[
\dim_k\left(
k[[x_1,\ldots,x_{n+1}]]\big/
  (f,\partial f/\partial x_i)
\right)=2.
\]
En coordenadas adecuadas $x_0,x_1,\ldots,x_{2m}$, la serie formal $f$
puede escribirse
\[
f=x_0^2+\sum_{i=1}^m x_ix_{i+m}+R,
\qquad\operatorname{ord}(R)\geq3.
\]
Sea $\mathfrak n'$ el ideal engendrado por $f$ y los
$\partial f/\partial x_i$, sea $\mathfrak n$ el ideal engendrado por
$x_0^2$ y los $x_i$ para $i\ne0$, y sea $\mathfrak q$ el ideal
maximal. Demostremos que $\mathfrak n=\mathfrak n'$. Claramente,
$\mathfrak n'\subset\mathfrak n$, y por Nakayama basta comprobar que
\[
\mathfrak n'+\mathfrak n\mathfrak q=\mathfrak n.
\]
Módulo $\mathfrak n\mathfrak q$, se tiene $f\equiv x_0^2$; módulo
$\mathfrak n\mathfrak q+(x_0^2)$, se tiene
\[
\frac{\partial f}{\partial x_i}\equiv x_{i+m},
\qquad
\frac{\partial f}{\partial x_{i+m}}\equiv x_i
\qquad(1\leq i\leq m),
\]
lo que demuestra la afirmación.

Bajo las hipótesis de 1.2.8, si $k=k(y)$, denotaremos por $a$ un
cuadrado cualquiera de $k$.

\pageoriginale

\subsubsection*{Lema 1.2.9.}
\underline{Sea $V$ un fibrado vectorial sobre un esquema $S$, y sea
$T$ un subesquema cerrado de $\mathbb A_S(V)$ localmente libre de
rango dos sobre $S$. Entonces $T$ está contenido en una y solo una
recta de $\mathbb A_S(V)$.}

La cuestión es local para la topología fppf sobre $S$, por lo que
podemos suponer que $T/S$ tiene una sección. Por traslación, podemos
suponer que $T$ contiene la sección cero. Sea $\mathfrak n$ el ideal
de $T$ que define $0$; es localmente libre de rango uno sobre $S$.
Toda forma lineal homogénea $v\in V$ sobre $\mathbb A_S(V)$ define un
elemento de $\mathfrak n$, y por tanto un epimorfismo
$V\twoheadrightarrow\mathfrak n$. Su núcleo define la recta buscada.

\subsubsection*{1.2.10.}
En particular, si $A$ es un subesquema cerrado radicial de rango dos
del espacio afín $\mathbb A_k^{n+1}$ sobre un cuerpo $k$ de
característica $2$, existe un $a\in k$ tal que, en coordenadas afines
adecuadas, $A$ tiene ecuaciones
\[
x_0^2-a=0,
\qquad x_i=0\quad(1\leq i\leq n).
\]
Basta tomar los $x_i$ que se anulan sobre la recta definida por $A$.

\subsubsection*{1.2.11.}
Demostremos 1.2.6 para $y$ degenerado ($p=2$ y $n=2m$). Podemos y
vamos a suponer que $k(y)=k(\sqrt a)$, que $Y$ está definido por una
ecuación $f$ en un esquema liso $Z$, y que existe un sistema de
coordenadas locales en $y$
\[
x=(x_i)_{0\leq i\leq2m}:
Z\longrightarrow\mathbb A_k^{2m+1}
\]
que lleva $y$ al punto cerrado $(\sqrt a,0,\ldots,0)$, tal que, en un
entorno de $y$,
\[
J^n_{Y/k}\simeq
\mathcal O_Y/(x_0^2-a,x_1,\ldots,x_{2m}).
\]
Tras reemplazar $f$ por un múltiplo adecuado, se tiene
\[
f=(x_0^2-a)+
 \sum_{1\leq i\leq j\leq2m}a_{ij}x_ix_j+R,
\]
donde
\begin{enumerate}[label=(\alph*),leftmargin=2.6em]
\item la forma cuadrática en $2m$ variables
$\sum a_{ij}x_ix_j$ es no degenerada;
\item después de extender los escalares a $\bar k$, $R$ tiene orden al
menos $3$ en $x_0-\sqrt a$ y los $x_i$.
\end{enumerate}

\pageoriginale

Sea $Y_0$ definido, como en 1.2.4, por
\[
Q=x_0^2-a+\sum a_{ij}x_ix_j.
\]

\underline{Sobre $Y$}, el sistema $(x_i)$ puede considerarse de nuevo
como una solución aproximada de $Q(X_i)=0$. Esta vez, el ideal $\delta$,
imagen inversa por $x$ del ideal engendrado por los
$\partial Q/\partial X_i$, satisface
$\mathcal O_Y/\delta\simeq J^n_{Y/k}$; en $y$ es
\[
\delta=(x_0^2-a,(x_i)_{1\leq i}).
\]
Sea $\mathfrak q$ el ideal maximal que define $y$. Para concluir como
en el caso no degenerado aplicando el teorema de la función implícita,
queda verificar que
\begin{equation}\tag{1.2.11.1}
Q(x_i)\equiv0\pmod{\delta^2\mathfrak q}
\qquad\underline{\text{sobre }Y}.
\end{equation}
Sea $\bar{\mathfrak q}$ el ideal que define $y$ después de extender
los escalares a una clausura algebraica de $k$. Se comprueba que
\[
\delta^2\bar{\mathfrak q}\cap\mathcal O_{Y,y}
=\delta^2\mathfrak q.
\]
Basta, por tanto, verificar \sisi{(4.11.1)}{(1.2.11.1)} tras extender
los escalares a $\bar k$; podemos suponer que $a=0$. En $y$ se tiene
entonces
\[
\delta^2\mathfrak q=
(x_0^5,(x_0^3x_i)_{1\leq i},
 (x_0x_ix_j)_{1\leq i,j},
 (x_ix_jx_k)_{1\leq i,j,k}).
\]
Se sigue que $\delta^2\mathfrak q\supset\mathfrak q^5$ y que
\[
x_0^2f\equiv x_0^4\pmod{\delta^2\mathfrak q},
\qquad
\delta^2\mathfrak q+(x_0^2f)\supset\mathfrak q^4,
\]
\[
x_0f\equiv x_0^3
 \pmod{\delta^2\mathfrak q+\mathfrak q^4},
\qquad
\delta^2\mathfrak q+(x_0f)\supset\mathfrak q^3.
\]

\pageoriginale

Por último,
\[
Q(x_i)\equiv0
\pmod{\delta^2\mathfrak q+(f)}
\qquad\underline{\text{sobre }Z},
\]
lo cual es (1.2.11.1).

\subsubsection*{Observación 1.2.12.}
Sea $Q$ una forma cuadrática no necesariamente homogénea, no nula,
en $n+1$ variables, y supongamos que el subesquema cerrado
$Y\subset\mathbb A_k^{n+1}$ de ecuación $Q$ tiene una singularidad
cuadrática ordinaria en un punto $y$. Un cambio lineal no homogéneo de
variables pone entonces $Y$ en la forma 1.2.3 o 1.2.4. En el caso no
degenerado, $k(y)=k$, y basta situar $y$ en el origen. En el caso
degenerado, basta que el soporte de $T^1_{Y/k}$ esté sobre la recta
$x_i=0$ $(1\leq i\leq2m)$ (cf. 1.2.9).

\subsection*{1.3. Módulos de singularidades cuadráticas}

La proposición siguiente resulta del teorema 1.1.4 de R. Elkik y de
cálculos explícitos que se dejan al lector (y que están esencialmente
contenidos en VI 6).

\subsubsection*{Proposición 1.3.1.}
\underline{Sea $S=\operatorname{Spec}(A)$ un esquema local henseliano
de punto cerrado $s$, y sea
$Y_0\subset\mathbb A_{k(s)}^{n+1}$ un esquema del tipo descrito en
1.2.3 o 1.2.4.}

\begin{enumerate}[label=(\roman*),leftmargin=2.6em]
\item \underline{Supongamos que $Y_0$ es del tipo descrito en 1.2.3,
definido por la ecuación}
\[
Q(\underline X)=\sum_{i\leq j}\bar a_{ij}x_i x_j=0
\qquad(\bar a_{ij}\in k(s)),
\]
\underline{y levantemos los $\bar a_{ij}$ a elementos $a_{ij}$ de
$A$. Entonces el $S$-esquema henseliano versal de deformaciones de la
singularidad de $Y_0$ en $y_0$ es}
\[
T=\operatorname{Spec}(A\{b\}),
\]
\underline{siendo la deformación versal el henselizado en $y_0$ del
$T$-esquema $Y\subset\mathbb A_T^{n+1}$ de ecuación}

\pageoriginale

\[
\sum a_{ij}x_i x_j-b=0.
\]

\item \underline{Supongamos que $Y_0$ es del tipo descrito en 1.2.4,
definido por la ecuación}
\[
Q(\underline x)=(x_0^2-\bar a)+\sum\bar a_{ij}x_i x_j=0
\qquad(\bar a,\bar a_{ij}\in k(s)),
\]
\underline{y levantemos los $\bar a_{ij}$ y $\bar a$ a elementos
$a_{ij}$ y $a$ de $A$. Entonces el $S$-esquema henseliano versal de
deformaciones de la singularidad de $Y_0$ en $y_0$ es}
\[
T=\operatorname{Spec}(A\{b,c\}),
\]
\underline{siendo la deformación versal el henselizado en $y_0$ del
$T$-esquema $Y\subset\mathbb A_T^{n+1}$ de ecuación}
\[
(x_0^2-a)+\sum a_{ij}x_i x_j+b x_0+c=0.
\]
\end{enumerate}

El corolario siguiente resulta inmediatamente de 1.2.6 y 1.3.1.

\subsubsection*{Corolario 1.3.2.}
\underline{Sea $S=\operatorname{Spec}(A)$ un esquema local henseliano
de punto cerrado $s$, y sea $f:X\to S$ un morfismo plano de
presentación finita. Supongamos que, en el punto cerrado $x$ de $X_s$,
el esquema $X_s$ tiene una singularidad cuadrática ordinaria, que
$k(x)$ es radicial sobre $k(s)$ y que $X_s$ tiene dimensión $n$.}

\begin{enumerate}[label=(\roman*),leftmargin=2.6em]
\item Si $x$ es no degenerado, \underline{existen una forma cuadrática
no degenerada $Q$ en $n+1$ variables con coeficientes en $A$, y un
elemento $b$ del ideal maximal de $A$, tales que el $S$-esquema
henselizado de $X$ en $x$ es isomorfo al henselizado en el origen del
subesquema cerrado de $\mathbb A_S^{n+1}$ de ecuación $Q-b=0$, como en
1.3.1(i).}

\item Si $x$ es degenerado
($n=2m$, $\operatorname{char}(k(s))=2$), \underline{existe una forma
cuadrática no homogénea en $n+1$ variables con coeficientes en $A$}
\begin{equation}\tag{\sisi{4.14.2}{1.3.2.1}}
Q(\underline X)=\sisi{x_0}{x_0^2}+b x_0+c+
 \sum_{1\leq i\leq j\leq 2m}a_{ij}x_i x_j,
\end{equation}

\pageoriginale

\underline{tal que $b$ pertenece al ideal maximal, la forma cuadrática
$\sum a_{ij}x_i x_j$ en $2m$ variables es no degenerada, y el
$S$-esquema henselizado de $X$ en $x$ es isomorfo al henselizado en
$(\sqrt c,0,\ldots,0)$ del subesquema cerrado de
$\mathbb A_S^{n+1}$ de ecuación $Q=0$.}
\end{enumerate}

\subsubsection*{Observación 1.3.3.}
En el caso (i), supongamos que el henselizado de $X_s$ en $x$ es
isomorfo, mediante un isomorfismo $\varphi$, al henselizado en $O$ del
subesquema cerrado $X_0\subset\mathbb A_k^{n+1}$ de ecuación
\[
\sum\bar a_{ij}x_i x_j=0,
\]
y supongamos que los $a_{ij}$ levantan los $\bar a_{ij}$. Se puede
tomar entonces
\[
Q=\sum a_{ij}x_i x_j.
\]

En el caso (ii), supongamos que el henselizado de $X_s$ en $x$ es
isomorfo, mediante un isomorfismo $\varphi$, al henselizado en
$(\sqrt a,0,\ldots,0)$ del subesquema cerrado
$X_0\subset\mathbb A_k^{n+1}$ de ecuación
\[
(x_0^2+\bar a)+\sum\bar a_{ij}x_i x_j=0,
\]
y supongamos que los $a_{ij}$ levantan los $\bar a_{ij}$. Se puede
tomar entonces
\[
Q=x_0^2+b x_0+c+\sum a_{ij}x_i x_j,
\qquad c\equiv\bar a\pmod{\mathfrak m_A}.
\]
En ambos casos, el isomorfismo prometido en 1.3.2 puede elegirse como
una prolongación de $\varphi$.

\subsubsection*{Corolario 1.3.4.}
\underline{Sea $f:X\to S$ como en 1.3.2. Existe un entorno $U$ de $x$
en $X$ tal que los puntos no lisos de $f$ contenidos en $U$ son puntos
cuadráticos ordinarios de sus fibras respectivas.}

\pageoriginale

\section*[2. Ciclos evanescentes: caso cuadrático ordinario]{2. Cálculo de los ciclos evanescentes en un caso cuadrático
ordinario estándar}

\subsection*{2.1. Preliminar: cohomología de un cono}

\subsubsection*{2.1.1.}
Sea $Y\subset\mathbb P^r$ una variedad proyectiva sobre un cuerpo
algebraicamente cerrado $k$. Utilizamos la notación siguiente.

El cono afín sobre $Y$ se denota por $X\subset\mathbb A^{r+1}$, y su
vértice por $O$;

$X_{(O)}$ es el henselizado de $X$ en $O$;

$X^*=X-\{O\}$ y $X^*_{(O)}=X_{(O)}-\{O\}$;

$X_1\subset\mathbb P^{r+1}$ es el cono proyectivo sobre $Y$, de modo
que $X=X_1-Y$;

$\widetilde X$, $\widetilde X_{(O)}$ y $\widetilde X_1$ son los
estallidos de $X$, $X_{(O)}$ y $X_1$ en $O$;

$Y_0$ es la imagen inversa de $O$ en $\widetilde X$,
$\widetilde X_{(O)}$ o $\widetilde X_1$;

$h:\widetilde X_1\to Y$ es la proyección natural;

$y_0,y_\infty:Y\to\widetilde X_1$ son las secciones de $h$ cuyas
imágenes son $Y_0$ e $Y$;

$F$ es un grupo abeliano de torsión de orden primo con la
característica de $k$. Todos los grupos de cohomología considerados a
continuación tienen coeficientes en $F$, que no se indicarán.

El objetivo de esta subsección es relacionar la cohomología de
$X^*_{(O)}$ con la de $Y$.

\subsubsection*{Proposición 2.1.2.}
\underline{Existen isomorfismos}
\[
\text{(i)}\qquad
H^i(X)\xrightarrow{\ \sim\ }H^i(\{O\})=
\begin{cases}
F,& i=0,\\
0,& i>0;
\end{cases}
\]
\[
\text{(ii)}\qquad
H^i_{\{O\}}(X)\xrightarrow{\ \sim\ }H^i_c(X).
\]
Esta proposición se demuestra mediante un argumento de homotopía, que
formalizamos ahora.

Sean $k$ un cuerpo algebraicamente cerrado, $U$ y $V$ dos $k$-esquemas,
$\Lambda$ un anillo de torsión de orden primo con la característica de
$k$, y
\[
K\in\operatorname{Ob}D^+(U,\Lambda),
\qquad L\in\operatorname{Ob}D^+(V,\Lambda).
\]

\pageoriginale

Por un ``morfismo'' $f$ de $(U,K)$ a $(V,L)$ entendemos un par formado
por un morfismo $f:U\to V$ y un morfismo $\varphi:f^*L\to K$. Tal
morfismo induce
\[
f^*:H^*(V,L)\longrightarrow H^*(U,K).
\]

Se dice que dos ``morfismos'' $f_0$ y $f_1$ son homotópicos si existen
un $k$-esquema conexo $T$ de tipo finito, dos puntos $0$ y $1$ de $T$,
y un ``morfismo'' $H$ de $(U\times T,\operatorname{pr}_1^*K)$ a
$(V,L)$ cuyas fibras en $0$ y $1$ inducen $f_0$ y $f_1$.

\subsubsection*{Lema 2.1.3.}
\underline{Si $f_0$ es homotópico a $f_1$, entonces
$f_0^*=f_1^*$.}

Existe una sucesión de puntos $0=x_0,x_1,\ldots,x_n=1$ de $T$,
curvas lisas conexas $\Gamma_i$, y morfismos $a_i:\Gamma_i\to T$
tales que $x_i$ y $x_{i+1}$ pertenecen a la imagen de $a_i$. Para
construir las $\Gamma_i$, normalicemos un subesquema cerrado conexo de
dimensión uno de $T$ que contenga $0$ y $1$. Esto nos reduce al caso en
que $T$ es una curva lisa. Consideremos el cuadrado conmutativo
\begin{center}
\begin{tikzcd}[row sep=large,column sep=large]
U\times T \arrow[r,"\operatorname{pr}_1"]
  \arrow[d,"\operatorname{pr}_2"']
  & U \arrow[d,"f"] \\
T \arrow[r,"t"'] & k .
\end{tikzcd}
\end{center}
Por cambio de base liso, aplicado a $t$, se tiene
\[
t^*Rf_*K\xrightarrow{\ \sim\ }
R\operatorname{pr}_{2*}(\operatorname{pr}_1^*K).
\]
La aplicación $f_i^*$ $(i=0,1)$ se factoriza, por tanto, por
\begin{center}
\begin{tikzcd}[row sep=large,column sep=large]
H^n(V,L) \arrow[r,"H^*"]
  & H^n(U\times T,\operatorname{pr}_1^*K) \arrow[r]
  & H^0\bigl(T,R^n\operatorname{pr}_{2*}
      (\operatorname{pr}_1^*K)\bigr)
      \arrow[dd,"\text{fibra en }i"] \\
{} & H^0\bigl(T,t^*H^n(U,K)\bigr) \arrow[ur,"\wr"] & {} \\
{} & {} & H^n(U,K) \arrow[ul,"\wr"']
\end{tikzcd}
\end{center}

\pageoriginale

y es, por consiguiente, independiente de $i$.

\subsubsection*{2.1.3.}
Demostremos 2.1.2. La multiplicación escalar proporciona una homotopía
entre la aplicación identidad de $X$ y la aplicación constante de $X$ a
$X$ de valor $O$; esto demuestra (i).

Utilizando multiplicaciones escalares cuyos factores tienden al
infinito, se ve que $Y$ es un retracto por deformación de $X_1^*$. La
afirmación (ii) resulta entonces del lema de los cinco y del diagrama
\begin{center}
\begin{tikzcd}[row sep=large,column sep=large]
{} \arrow[r,"\partial"]
  & H^i_{\{O\}}(X) \arrow[r] \arrow[d]
  & H^i(X_1) \arrow[r] \arrow[d,equal]
  & H^i(X_1^*) \arrow[r,"\partial"] \arrow[d,"\wr"']
  & {} \\
{} \arrow[r,"\partial"]
  & H^i_c(X) \arrow[r]
  & H^i(X_1) \arrow[r]
  & H^i(Y) \arrow[r,"\partial"]
  & {} .
\end{tikzcd}
\end{center}

\subsubsection*{Corolario 2.1.4.}
\[
H^i(X^*)\xrightarrow{\ \sim\ }H^i(X^*_{(O)}).
\]
En efecto, apliquemos el lema de los cinco a
\begin{center}
\begin{tikzcd}[row sep=large,column sep=large]
{} \arrow[r]
  & H^i_{\{O\}}(X) \arrow[r] \arrow[d,"\wr"]
  & H^i(X) \arrow[r]
      \arrow[d,"\wr"',"{(2.1.2(i))}"]
  & H^i(X^*) \arrow[r] \arrow[d]
  & {} \\
{} \arrow[r]
  & H^i_{\{O\}}(X_{(O)}) \arrow[r]
  & H^i(X_{(O)}) \arrow[r]
  & H^i(X^*_{(O)}) \arrow[r]
  & {} .
\end{tikzcd}
\end{center}

\subsubsection*{2.1.5.}
Las sucesiones espectrales de Leray para las restricciones de $h$ a
$\widetilde X_1-Y_0$ y $\widetilde X_1-Y$ muestran que las aplicaciones
de restricción
\[
H^i(\widetilde X_1-Y_0)\xrightarrow{\ \sim\ }H^i(Y),
\qquad
H^i(\widetilde X)\xrightarrow{\ \sim\ }H^i(Y_0)=H^i(Y)
\]
son isomorfismos. Bajo estos isomorfismos, las sucesiones exactas largas
de cohomología de los pares $(\widetilde X_1-Y_0,Y)$ y
$(\widetilde X_1-Y,Y_0)$ forman las filas del diagrama

\pageoriginale

\begin{equation}\tag{2.1.5.1}
\begin{tikzcd}[column sep=large]
H^{i-1}(X^*) \arrow[r] \arrow[d,equal]
  & H^{i-2}(Y)(-1) \arrow[r,"\eta"'] \arrow[d,equal]
  & H^i(Y) \arrow[r] \arrow[d,dash,"-1"]
  & H^i(X^*) \arrow[r] \arrow[d,equal]
  & {} \\
H^{i-1}(X^*) \arrow[r]
  & H^{i-2}(Y)(-1) \arrow[r,"-\eta"']
  & H^i(Y) \arrow[r]
  & H^i(X^*) \arrow[r]
  & {} .
\end{tikzcd}
\end{equation}

\subsubsection*{Lema 2.1.6.}
\underline{Las flechas denotadas por $\eta$ y $-\eta$ en (2.1.5.1)
son productos cup con la clase de cohomología de una sección
hiperplana. El diagrama (2.1.5.1) es conmutativo.}

La primera afirmación resulta de que $\eta$ (respectivamente,
$-\eta$) es la restricción a $Y$ (respectivamente, a $Y_0\simeq Y$)
de la clase de cohomología del haz invertible $\mathcal O(Y)$
(respectivamente, $\mathcal O(Y_0)$) sobre $\widetilde X_1-Y_0$
(respectivamente, $\widetilde X_1-Y$).

Utilicemos la sucesión espectral de Leray para $h$. Las dos filas de
(2.1.5.1) se obtienen aplicando el funtor $H(Y,-)$ a los triángulos
distinguidos
\begin{center}
\small
\begin{tikzcd}[column sep=normal]
{} & y_{\infty*}Ry_\infty^!F \arrow[r]
  & R(h|\widetilde X_1-Y_0)_*F
      \arrow[r,"\text{\textcircled{\scriptsize 1}}"]
  & R(h|X^*)_*F \arrow[r,"\text{\textcircled{\scriptsize 2}}"]
  & {} \\
\text{y}
  & y_{0*}Ry_0^!F \arrow[r]
  & R(h|\widetilde X_1-Y)_*F \arrow[r]
  & R(h|X^*)_*F \arrow[r]
  & {} .
\end{tikzcd}
\end{center}
Los haces de cohomología no nulos de los complejos en los vértices de
estos triángulos aparecen en los grados $(2)$, $(0)$ y $(0$ y $1)$.
Por tanto, las flechas \textcircled{\scriptsize 1} y
\textcircled{\scriptsize 2} están determinadas por su efecto sobre los
haces de cohomología. Esto nos reduce al caso en que $Y$ es un punto,
que se deja al lector.

\subsubsection*{2.1.7.}
La sucesión exacta larga del par $(\widetilde X_{(O)},Y_0)$ es un
análogo local de la segunda fila de (2.1.5.1). Se tiene
\[
H^i(\widetilde X_{(O)})\xrightarrow{\ \sim\ }H^i(Y_0)
\]
por cambio de base propio para el morfismo
$\widetilde X_{(O)}\to X_{(O)}$. La sucesión exacta larga en cuestión
puede escribirse

\pageoriginale

\begin{equation}\tag{\sisi{2.1 7.1}{2.1.7.1}}
\begin{tikzcd}[column sep=large]
H^{i-1}(X^*_{(O)}) \arrow[r]
  & H^{i-2}(Y_0)(-1) \arrow[r,"-\eta"]
  & H^i(Y_0) \arrow[r]
  & H^i(X^*_{(O)}) & {} \\
H^{i-1}(X^*) \arrow[r] \arrow[u,"\wr"]
  & H^{i-2}(Y)(-1) \arrow[r,"-\eta"] \arrow[u,equal]
  & H^i(Y) \arrow[r] \arrow[u,equal]
  & H^i(X^*) \arrow[u,dash,"\wr"]
  & {} ,
\end{tikzcd}
\end{equation}
siendo la segunda fila de este diagrama conmutativo la segunda fila de
(2.1.5.1).

El lema técnico siguiente se utilizará en 2.2.7.

\ifthenelse{\boolean{original}}
  {\subsubsection*{Lema 2.7.8.}}
  {\subsubsection*{Lema 2.1.8.}}
\begin{center}
\footnotesize
\begin{tikzcd}[column sep=normal]
H^{n-1}(Y_0) \arrow[d,equal]
  & H^{n-1}(\widetilde X) \arrow[l,"\sim"'] \arrow[r]
  & H^{n-1}(X^*) \arrow[r]
  & H^n_{\{x\}}(X) \arrow[r]
  & H^n_c(X_s) \arrow[d,equal] \\
H^{n-1}(Y) \arrow[rrrr,"\partial"]
  & {} & {} & {} & H^n_c(X_s)
\end{tikzcd}
\end{center}
\underline{es anticonmutativo.}

Por 2.1.7, esto equivale a demostrar la conmutatividad de
\begin{center}
\begin{tikzcd}[column sep=large,row sep=large]
H^{n-1}(X_1-\{O\}) \arrow[r,"\partial"] \arrow[d,"\wr"]
  & H^n_{\{O\}}(X_1)=H^n_{\{O\}}(X) \arrow[d] \\
H^{n-1}(Y) \arrow[r,"\partial"]
  & H^n_c(X).
\end{tikzcd}
\end{center}
Esto es formal.

\subsection*{2.2. Ciclos evanescentes}

\subsubsection*{2.2.1.}
Sea $S$ un trait henseliano, y sean $s,\eta,\bar s,\bar\eta$ como en
XIII, 0.2.5. Pongamos $\Lambda=\mathbb Z/k$, donde $k$ es primo con
la característica residual de $S$. Sea $\mathbb A^{n+1}_S$ el espacio
afín estándar de dimensión $n+1$ sobre $S$, y sea $X$ el subesquema
cerrado de $\mathbb A^{n+1}_S$ definido por una ecuación cuadrática
$Q=0$ no nula módulo un uniformizante:

\pageoriginale

\[
Q(\underline x)=\sum_{i\leq j}a_{ij}x_ix_j+\sum_i b_ix_i+c.
\]
Consideremos $\mathbb A^{n+1}_S$ como el complemento en
$\mathbb P^{n+1}_S$ del hiperplano del infinito $H$. Entonces $X$ se
obtiene de la cuádrica $X_1\subset\mathbb P^{n+1}_S$ de ecuación
homogénea
\begin{equation}\tag{2.2.1.1}
\sum_{i\leq j}a_{ij}x_ix_j+\sum_i b_ix_iz+cz^2=0
\end{equation}
quitando $Y=X_1\cap H$.

Hacemos las hipótesis siguientes:

\begin{enumerate}
\item[(a)] $Y$ es una cuádrica lisa sobre $S$, es decir, la forma
cuadrática $\sum_{i\leq j}a_{ij}x_ix_j$ es ordinaria;

\item[(b)] la cuádrica $\sisi{X'_s}{X_s}$ es singular, es decir,
$X_{\bar s}$ es un cono cuadrático.
\end{enumerate}

Denotemos por $x_0$ el punto singular de $X_s$, y sea $f:X\to S$ la
proyección.

\subsubsection*{Observación 2.2.2.}
Si $x_0$ es un punto racional, una traslación lo lleva al origen de las
coordenadas. Si $x_0=0$, los $b_i$ y $c$ son divisibles por un
uniformizante.

El subesquema $A$ de $X_s$ definido por las ecuaciones
$\partial Q/\partial x_i$ está concentrado en $x_0$. Salvo en el caso
excepcional en que $k(s)$ tiene característica $2$ y $n+1$ es par, este
esquema tiene grado uno sobre $k(s)$ y, por tanto, $x_0$ es racional.
En el caso excepcional es ramificado sobre $k(s)$ y de rango $2$, de
modo que $k(x_0)$ es o bien $k(s)$, o bien una extensión cuadrática
puramente inseparable de $k(s)$. Las afirmaciones sobre $A$ se comprueban
más fácilmente tras pasar a la clausura algebraica $k(\bar s)$ de
$k(s)$.

\subsubsection*{Proposición 2.2.3.}
\underline{Bajo las hipótesis de 2.2.1, los morfismos}

\pageoriginale

\[
H^i(X_{\bar\eta},\Lambda)
 \xrightarrow{\ \sim\ }
\mathbb H^i\bigl(X_{\bar s},R\psi_{\bar\eta}(\Lambda)\bigr)
 \xrightarrow{\ \sim\ }
R^i\psi_{\bar\eta}(\Lambda)_{x_0}
\]
\[
H_c^i(X_{\bar\eta},\Lambda)
 \xleftarrow{\ \sim\ }
\mathbb H_c^i\bigl(X_{\bar s},R\psi_{\bar\eta}(\Lambda)\bigr)
 \xleftarrow{\ \sim\ }
\mathbb H^i_{\{x_0\}}\bigl(X_{\bar s},R\psi_{\bar\eta}(\Lambda)\bigr)
\]
\underline{son isomorfismos.}

Que las flechas de la izquierda sean isomorfismos resulta de 2.1.8.6
y 2.1.10.5; la afirmación para las flechas de la derecha es un
corolario de 2.1.2. En efecto, existe un triángulo distinguido
\[
(\Lambda\text{ sobre }X_{\bar s})[0]
 \longrightarrow R\psi_{\bar\eta}(\Lambda)
 \longrightarrow R\Phi_{\bar\eta}(\Lambda),
\]
y 2.1.2 se aplica a $(\Lambda\text{ sobre }X_{\bar s})$, mientras que
$R\Phi_{\bar\eta}(\Lambda)$ está soportado en $x_0$.

\subsubsection*{Corolario 2.2.4.}
\underline{Si la fibra genérica geométrica $X_{\bar\eta}$ es singular,
es decir, si sigue siendo un cono cuadrático, entonces todos los haces
$R^i\Phi(\Lambda)$ son nulos.}

Por 2.1.2,
\[
\Lambda=H^{\sisi{o}{0}}(X_{\bar s},\Lambda)
 \xrightarrow{\ \sim\ }H^{\sisi{o}{0}}(X_{\bar\eta},\Lambda),
\qquad
0=H^i(X_{\bar s},\Lambda)=H^i(X_{\bar\eta},\Lambda)
\quad(i>0).
\]
Se aplica entonces la sucesión exacta larga XIII, 2.1.8.9 (válida por
XIII, 2.1.9).

\subsubsection*{2.2.5.}
Supongamos ahora que
\[
(*)\qquad\text{la fibra genérica $X_\eta$ es lisa.}
\]
Para simplificar, supondremos también que $S$ es estrictamente
henseliano. Los haces de ciclos evanescentes se calculan fácilmente con
ayuda de 2.2.3 y XII, 3.7. La cohomología con coeficientes en $\Lambda$
considerada aquí resulta de la cohomología $\ell$-ádica del pasaje
citado mediante la fórmula de coeficientes universales (cf. 2.1.13).
Los resultados son los siguientes.

\medskip
\noindent\underline{A. Cohomología sin soportes.}

Se tiene $R^i\psi_{\sisi{\eta}{\bar\eta}}(\Lambda)=0$ para
$i\ne0,n$.

\pageoriginale

\begin{enumerate}
\item[1)] \underline{Si $n\ne0$:}
$\psi_{\bar\eta}(\Lambda)=\Lambda$. Además,
$R^n\psi_{\bar\eta}(\Lambda)$ es isomorfo (no canónicamente) al haz
$\Lambda$ sobre $x_0$, prolongado por cero a $X_s$.

\item[2)] \underline{Para $n$ arbitrario:}
incluso cuando $n=0$, se comprueba que
$R^i\Phi_{\bar\eta}(\Lambda)=0$ para $i\ne n$, y que
$R^n\Phi_{\bar\eta}(\Lambda)$ es isomorfo (no canónicamente) al haz
$\Lambda$ sobre $x_0$, prolongado por cero a $X_s$.
\end{enumerate}

\medskip
\noindent\underline{B. Cohomología con soporte en $x_0$.}

Se tiene
\[
H^i_{\{x_0\}}\bigl(X_s,R\psi_{\bar\eta}(\Lambda)\bigr)=0
\qquad(i\ne n,2n).
\]
Se dispone, además, de un morfismo traza
\begin{equation}\tag{2.2.5.1}
\operatorname{tr}:
H^{2n}_{\{x_0\}}\bigl(X_s,R\psi_{\bar\eta}(\Lambda(n))\bigr)
\longrightarrow\Lambda.
\end{equation}
\underline{Si $n\ne0$:} el morfismo (2.2.5.1) es un isomorfismo.
Además,
\[
H^n_{\{x_0\}}\bigl(X_s,R\psi_{\bar\eta}(\Lambda(n))\bigr)
\]
es isomorfo (no canónicamente) a $\Lambda$.

\medskip
\noindent\underline{C. Dualidad.}
Para
\[
a\in R^n\psi_{\bar\eta}(\Lambda)_{x_0},
\qquad
b\in H^n_{\{x_0\}}\bigl(X_s,R\psi_{\bar\eta}(\Lambda(n))\bigr),
\]
pongamos $(a,b)=\operatorname{Tr}(a\wedge b)$, donde
$\operatorname{Tr}$ es el morfismo XIII, 2.2.6
\[
\operatorname{Tr}:
H^{2n}_{\{x_0\}}\bigl(X_s,R\psi_{\sisi{\eta}{\bar\eta}}(\Lambda(n))\bigr)
\longrightarrow\Lambda.
\]
Resulta de la dualidad de Poincaré sobre $X_{\bar\eta}$ y de los
isomorfismos 2.2.3 que $(a,b)$ proporciona una dualidad perfecta entre
los $\Lambda$-módulos libres
\[
R^n\psi_{\bar\eta}(\Lambda)_{x_0}
\quad\text{y}\quad
H^n_{\{x_0\}}\bigl(X_s,R\psi_{\bar\eta}(\Lambda(n))\bigr).
\]

\medskip
\noindent\underline{D. $n=2m>0$.}

Por XII, 3.7, los grupos
\[
H^n_{\{x_0\}}\bigl(X_s,R\psi_{\bar\eta}(\Lambda(m))\bigr)
\quad\text{y}\quad
R^n\psi_{\bar\eta}(\Lambda(m))_{x_0}
\]
tienen generadores naturales, definidos salvo signo; denotémoslos por
$\delta$ y $\delta'$, respectivamente. Sea $\phi$ la aplicación natural
del primero de estos grupos al segundo. Los signos de $\delta$ y
$\delta'$ pueden normalizarse de modo que
\begin{equation}\tag{2.2.5.2}
(\delta',\delta)=1;
\end{equation}
entonces

\pageoriginale

\begin{equation}\tag{2.2.5.3}
\phi(\delta)=(-1)^m\,2\,\delta'.
\end{equation}

Sea $Z$ la extensión cuadrática separable de $k(\eta)$ que es el centro
de la parte par del álgebra de Clifford de la forma cuadrática
(2.2.1.1), y sea $\epsilon$ el carácter correspondiente, de orden $1$
o $2$, del grupo de inercia $I$. Para $\sigma\in I$, se tiene
\begin{equation}\tag{2.2.5.4}
\sigma\delta=\epsilon(\sigma)\delta
\end{equation}
y
\begin{equation}\tag{2.2.5.5}
\sigma\delta'=\epsilon(\sigma)\delta'.
\end{equation}
Las fórmulas (2.2.5.2), (2.2.5.3) y (2.2.5.4), en cohomología con
coeficientes en $\mathbb Z/2k$, implican que la variación es
\begin{equation}\tag{2.2.5.6}
\operatorname{Var}(\sigma)(a)=
\frac{\epsilon(\sigma)-1}{2}\,(-1)^m\,
\sisi{(a\delta)}{(a,\delta)}\,\delta.
\end{equation}

\medskip
\noindent\underline{E. $n=2m+1$---Ciclos evanescentes.}

Por XII, 3.7, los grupos
\[
H^n_{\{x_0\}}\bigl(X_s,R\psi_{\bar\eta}(\Lambda(m))\bigr)
\quad\text{y}\quad
R^n\psi_{\bar\eta}(\Lambda(m+1))_{x_0}
\]
tienen generadores naturales, definidos salvo signo; denotémoslos por
$\delta$ y $\delta'$, respectivamente. Con la notación de D, sus signos
pueden normalizarse de modo que
\begin{equation}\tag{2.2.5.7}
(\delta',\delta)=1;
\end{equation}
entonces
\begin{equation}\tag{2.2.5.8}
\phi(\delta)=0.
\end{equation}

\medskip
\noindent\underline{F. $n=2m+1$---La variación.}

El grupo de inercia $I$ actúa trivialmente sobre la cohomología de
$X_\eta$. Esto resulta de XII, 3.7 y del hecho de que $I$ actúa
trivialmente sobre la cohomología de $Y_\eta$, pues $Y$ es propio y
liso sobre $S$.

Para $\sigma$ en el grupo de inercia se tiene necesariamente
\begin{equation}\tag{2.2.5.9}
D(\sigma)(\delta')=\lambda(\sigma)\delta,
\qquad \lambda(\sigma)\in\Lambda(1).
\end{equation}
La fórmula de aditividad (XIII, 1.4.3.4) se convierte en
\[
\lambda(\sigma\tau)=\lambda(\sigma)+\lambda(\tau).
\]
Sea
\[
\epsilon:I\longrightarrow\Lambda(1)=\mu_k
\]
el homomorfismo tal que

\pageoriginale

\[
\sigma\bigl(\sqrt[k]{t}\bigr)
=\epsilon(\sigma)\sqrt[k]{t}
\]
para un uniformizante $t$. En vista de la estructura del grupo de
inercia moderada, $\lambda$ es necesariamente un múltiplo de $\epsilon$.
La fórmula (2.2.5.9) puede reescribirse
\begin{equation}\tag{2.2.5.10}
\operatorname{Var}(\sigma)(a)
=\lambda_X\epsilon(\sigma)\sisi{(a\delta)}{(a,\delta)}\delta
\end{equation}
para cierto $\lambda_X\in\Lambda$ (que depende de $X/S$).
Determinaremos $\lambda_X$ más adelante.

\subsubsection*{2.2.6.}
Daremos una descripción local para la topología étale de los ciclos
$\delta$ y $\delta'$ de D y E.

El caso de $n$ par es fácil. Si $2\nmid k$
($\Lambda=\mathbb Z/k$), entonces $\pm\delta$ está caracterizado por
(2.2.5.3) y (2.2.5.4), es decir, por
\begin{equation}\tag{2.2.6.1}
(\delta,\delta)=(-1)^m\,2.
\end{equation}
En general, $\pm\delta$ es la imagen de cualquier ciclo
$\widetilde\delta$, en cohomología con coeficientes en
$\mathbb Z/2^ak$ para $a$ suficientemente grande, que satisface
(2.2.6.1).

Supongamos que $n$ es impar. Entonces $x_0$ es un punto racional
(2.2.2). Sea $X_{s(0)}$ el henselizado de $X_s$ en $x_0$, sea
$\widetilde X_{s(0)}$ su estallido en $x_0$, sea $Y_0$ la imagen
inversa de $x_0$ en $\widetilde X_{s(0)}$ (una cuádrica), y pongamos
\[
X^*_{s(0)}=X_{s(0)}-\{x_0\}
=\widetilde X_{s(0)}-Y_0.
\]
Consideremos la aplicación compuesta (cf. XIII, 2.4.6.1)
\begingroup
\small
\begin{equation}\tag{2.2.6.2}
\begin{tikzcd}[column sep=large,row sep=large]
H^{n-1}(Y_0,\Lambda(m))
  & H^{n-1}(\widetilde X_{s(0)},\Lambda(m))
      \arrow[l,"\sim"'] \arrow[r]
  & H^{n-1}(X^*_{s(0)},\Lambda(m)) \arrow[d] \\
{}
  & H^n_{\{\sisi{x}{x_0}\}}\bigl(X_s,R\Psi_{\bar\eta}(\Lambda(m))\bigr)
  & H^n_{\{x_0\}}(X_s,\Lambda(m)) \arrow[l]
\end{tikzcd}
\end{equation}
\endgroup

\subsubsection*{Lema 2.2.7.}
\underline{La aplicación (2.2.6.2) identifica
$H^n_{\{\sisi{x}{x_0}\}}(X_s,R\Psi_{\bar\eta}(\Lambda(m)))$ con el
cociente primitivo de la cohomología de dimensión media de la cuádrica
$Y_0$ de dimensión $2m$; las dos clases $\pm\delta$}

\pageoriginale

\underline{son las imágenes de los generadores naturales de este
cociente primitivo.}

Los ciclos $\pm\delta$ están caracterizados por sus imágenes en
$H_c^{\sisi{i}{n}}(X_{\bar\eta},\Lambda(m))$. Utilizando (2.1.7.1),
basta, por tanto, demostrar que la aplicación compuesta
\begin{center}
\begin{tikzcd}[column sep=huge,row sep=huge]
H^{n-1}(Y_0,\Lambda(m))
  & H^{n-1}(\widetilde X_s,\Lambda(m))
      \arrow[l,"\sim"'] \arrow[r]
  & H^{n-1}(X_s^*,\Lambda(m)) \arrow[d] \\
H_c^n(X_{\bar\eta},\Lambda(m))
  & H_c^n(X_s,\Lambda(m)) \arrow[l,"\operatorname{sp}"']
  & H^n_{\{x_0\}}(X_s,\Lambda(m)) \arrow[l]
\end{tikzcd}
\end{center}
envía los generadores distinguidos del cociente primitivo del primer
grupo a los generadores distinguidos del último grupo.

Apliquemos 2.1.8. Salvo signo, la aplicación puede reescribirse
\begin{center}
\begin{tikzcd}[column sep=huge,row sep=huge]
H^{n-1}(Y_s,\Lambda(m)) \arrow[r,"\partial"] \arrow[d,equal]
  & H_c^n(X_s,\Lambda(m)) \arrow[d] \\
H^{n-1}(Y_{\bar\eta},\Lambda(m)) \arrow[r]
  & H_c^n(X_{\bar\eta},\Lambda(m)),
\end{tikzcd}
\end{center}
y de ello resulta 2.2.7.

\section*{3. La fórmula de Picard--Lefschetz}

\subsection*{3.1. Resultados preliminares}

\subsubsection*{3.1.1.}
Sea $S$ un trait henseliano y sea $f:X\to S$ un morfismo plano de tipo
finito, puramente de dimensión relativa $n$. Supongamos que la fibra
especial $X_s$ es lisa salvo en un conjunto finito de puntos $\Sigma'$,
en los cuales $X_s$ tiene una singularidad cuadrática ordinaria.

\pageoriginale

Sea $\Sigma$ el conjunto de puntos $x\in\Sigma'$ tales que la fibra
genérica $X_\eta$ es lisa en un entorno de $x$. Por último, sea
$\Lambda$ un anillo de torsión cuyo orden es primo con la característica
residual de $S$.

\subsubsection*{Proposición 3.1.2.}
\begin{enumerate}
\item[(i)] \underline{Los haces $R^i\Phi(\Lambda)$ son nulos para
$i\ne n$.}

\item[(ii)] \underline{El haz $R^n\Phi(\Lambda)$ es
\sisi{el único}{nulo} fuera de $\Sigma$. Su restricción a $\Sigma$ es
un haz de $\Lambda$-módulos de rango uno.}

\item[(iii)] \underline{Para todo punto geométrico $\bar x$ de
$\Sigma$, el emparejamiento $(a,b)$ de 2.2.5(C) pone en dualidad los
$\Lambda$-módulos}
\[
\bigl(R^n\psi(\Lambda)\bigr)_{\bar x}
\quad\text{y}\quad
H^n_{\{\bar x\}}\bigl(
\sisi{R^n\psi_{\bar\eta}(\Lambda)}
{X_{\bar s},R\psi_{\bar\eta}(\Lambda)}\bigr)
\]
\underline{. Son libres de rango uno si $n\ne0$, y de rango dos si
$n=0$.}
\end{enumerate}

Las afirmaciones de 3.1.2 son locales para la topología étale. Por
tanto, basta tratar el caso en que $S$ es estrictamente henseliano y
$X$ es la hipersuperficie del espacio afín relativo
$\mathbb A^{n+1}_S$ definida por una ecuación cuadrática que satisface
2.2.1(a) y (b) (aplicando 1.3.2).

El morfismo $f$ es liso en todo punto de $X_s-\Sigma'$, por lo que
$R\Phi=0$ fuera de $\Sigma'$. Por 2.2.4, $R\Phi$ se anula incluso fuera
de $\Sigma$. Podemos suponer ahora 2.2.5(*); las afirmaciones (ii) y
(iii) resultan entonces de 2.2.5(A), (B) y (C).

\subsection*{3.2. Dimensión relativa par}

Conservamos la notación de 3.1 y tomamos $n=2m$ par. Para simplificar,
suponemos que $S$ es estrictamente henseliano (el caso general se
obtendría por descenso).

\subsubsection*{Proposición 3.2.1.}
\begin{enumerate}
\item[(i)] \underline{Para $x\in\Sigma$, el grupo
$H^n_{\{x\}}\bigl(
\sisi{R^n\psi_{\bar\eta}(\Lambda(m))}
{X_s,R\psi_{\bar\eta}(\Lambda(m))}\bigr)$ tiene un generador natural
$\delta$, bien definido salvo signo, caracterizado por ser natural en
$\Lambda$ y por satisfacer}

\pageoriginale

\[
(\delta,\delta)=(-1)^m\,2.
\]
\underline{Para $n=0$, añádase la condición
$\operatorname{Tr}(\delta)=0$.}

\item[(ii)] \underline{Para un carácter adecuado, independiente de
$\Lambda$,}
\[
\varepsilon_x:I\longrightarrow\{\pm1\}
\]
\underline{del grupo de inercia $I$, la variación es}
\[
\operatorname{Var}(\sigma)(a)
=(-1)^m\left(\frac{\varepsilon_x(\sigma)-1}{2}\right)
\sisi{(a\delta)}{(a,\delta)}\delta
\]
\underline{para $a\in R^n\Phi_{\bar\eta}(\Lambda(m))$. En
consecuencia,}
\[
\sigma(\delta)=\varepsilon_x(\sigma)\delta.
\]
\end{enumerate}

Pasando al límite inverso, basta tratar el caso universal
$\Lambda=\mathbb Z_\ell$. Salvo signo, solo un $\delta$ puede entonces
satisfacer las condiciones de (i). Como en 3.1.2, nos reducimos al caso
2.2.5(*) y aplicamos 2.2.5(D).

\subsubsection*{Complemento 3.2.2.}
Supongamos que, en $x\in\Sigma$, el henselizado $X_{(x)}$ de $X$ en
$x$ es isomorfo al henselizado en $x_0$ del esquema proyectivo $Q$ de
ecuación
\[
\sum_{i\leq j}a_{ij}x_ix_j=0.
\]
Entonces el carácter $\varepsilon_x$ está definido por la extensión
cuadrática separable $Z(C^+(Q))$ del cuerpo de fracciones. Esto resulta
de 2.2.5(D).

\subsubsection*{3.2.3.}
En característica residual distinta de $2$, el Complemento 3.2.2 puede
hacerse más explícito. En este caso, $I$ tiene un único carácter no
trivial $\varepsilon$ de orden $2$. Para un uniformizante $t$, se tiene
\[
\sigma(\sqrt t)=\varepsilon(\sigma)\sqrt t.
\]
Por otra parte, $X_{(x)}$ es entonces isomorfo al henselizado en $0$ de
un esquema afín

\pageoriginale

\[
\sum a_{ij}x_ix_j=b,
\]
donde $b$ pertenece al ideal maximal y
$Q(x)=\sum a_{ij}x_ix_j$ es no degenerada. El ideal $(b)$ define la
imagen del subesquema de $X_{(x)}$ que es el soporte de
$J^n_{X_{(x)}/S}$. El centro del álgebra de Clifford es isomorfo a
\[
k(\eta)\left(\left((-1)^{m+1}2b\det(a_{ij})\right)^{1/2}\right)
\]
(Bourbaki, \emph{Alg\`ebre}, cap. 9, \S9,
n\textsuperscript{o}~4, observación final), de donde
\[
\varepsilon_x=\varepsilon^{v(b)}.
\]
La variación es nula si $b$ tiene valuación par; en caso contrario,
$\varepsilon_x=\varepsilon$.

\subsection*{3.3. Dimensión relativa impar}

Para demostrar la fórmula de Picard--Lefschetz en dimensión relativa
impar, tendremos que reducirnos al caso trascendente tratado en XIV,
3.2.1.1. El argumento de especialización necesario nos obliga a trabajar
sobre una base que ya no es unidimensional. Utilizaremos argumentos
\emph{ad hoc} adaptados al caso preciso que necesitamos: no dispongo de
una teoría utilizable sobre una base de dimensión mayor que uno, y he
preferido no cubrir mi ignorancia con generalidades.

\subsubsection*{3.3.1.}
Sea $S=\operatorname{Spec}(A)$ un esquema estrictamente local
noetheriano, y sea $f:X\to S$ un morfismo plano de tipo finito,
puramente de dimensión relativa $n$, cuya fibra especial $X_s$ no tiene
más puntos no lisos que puntos cuadráticos no degenerados. Sea $s$ el
punto cerrado de $S$, sea $X_s$ la fibra especial y sea
$x_0\in X_s$ un punto singular.

Por 1.3.2(i), en un entorno de $x_0$ el esquema $X/S$ es localmente
isomorfo a un $S$-esquema definido en $\mathbb A^{n+1}_S$ por una
ecuación $Q-b=0$, donde $b$ pertenece al ideal maximal y $Q$ es una
forma cuadrática no degenerada. El ideal $(b)$ de $A$ solo depende de
$(X/S,x_0)$: en un entorno de $x_0$ existe una sección
\[
s\in X(A/(b))
\]
de $X$ sobre $\operatorname{Spec}(A/(b))$ tal que
\[
J^n_{X/S}\simeq
\mathcal O_{s(\operatorname{Spec}(A/(b)))}.
\]
Denotamos esta sección por $s_x$, y por $b(x)$ un generador cualquiera
de $(b)$.

\pageoriginale

\subsubsection*{3.3.2.}
Supongamos que $S$ es normal, y sea $b\in A$. Si $n$ es invertible en
$A$ y $b\ne0$, denotemos por $\varepsilon_b^{(n)}$ el homomorfismo
siguiente del grupo de Galois del cuerpo de fracciones de $A$ en
$\mathbb Z/(n)(1)$:
\[
\varepsilon_b^{(n)}(\sigma)
= (\sigma\sqrt[n]{b})/\sqrt[n]{b}.
\]
Se tiene
\[
\bigl(\varepsilon_b^{(nm)}\bigr)^m=\varepsilon_b^{(n)}.
\]
Si $\Lambda$ es una $\mathbb Z/n$-álgebra, denotamos de nuevo por
$\varepsilon_b$ la composición
\[
\operatorname{Gal}\longrightarrow\mathbb Z/(n)(1)
\longrightarrow\Lambda(1).
\]
Es independiente de $n$.

El carácter $\varepsilon_b$ se factoriza por
\[
\pi_1\bigl(S-(\text{el subesquema }b=0)\bigr).
\]

\subsubsection*{3.3.3.}
Supongamos que $S$ es un trait estrictamente henseliano y que $n$ es
impar, $n=2m+1$. Sea $x$ un punto singular de $X_s$ en un entorno del
cual $X_\eta$ es liso, y sea $X_1/S$ un subesquema de
$\mathbb A^{n+1}_S$ de ecuación $Q-b=0$, como en 1.3.2(i). La
hipótesis sobre $x$ significa que $b$, aunque nulo módulo el ideal
maximal, no es nulo.

Sea $X_{s(x)}$ el henselizado de $X_s$ en $x$, sea
$\widetilde X_{s(x)}$ el estallido de $X_s$ en $x$, sea $Y$ la
cuádrica imagen inversa de $x$, y pongamos
\[
\sisi{\widetilde X^*_{s(x)}}{X^*_{s(x)}}
=X_{s(x)}-\{x\}=\widetilde X_{s(x)}-Y.
\]
El morfismo compuesto
\begin{center}
\begin{tikzcd}[column sep=huge,row sep=huge]
H^{n-1}(Y,\Lambda(m))
  & H^{n-1}(\widetilde X_{s(x)},\Lambda(m))
      \arrow[l,"\sim"'] \arrow[r]
  & H^{n-1}(X^*_{s(x)},\Lambda(m)) \arrow[d] \\
{}
  & H^n_{\{x\}}\bigl(X_s,R\Psi_{\bar\eta}(\Lambda(m))\bigr)
  & H^n_{\{x\}}(X_s,\Lambda(m)) \arrow[l]
\end{tikzcd}
\end{center}
identifica el último grupo con el cociente primitivo (XII, 3.5) de la
cohomología de la cuádrica $Y$. En efecto, esto resulta de 2.2.7 en
una situación localmente isomorfa. Las imágenes de los dos generadores
naturales de este cociente (las clases de cohomología de las familias de
generatrices de $Y$) son los dos vectores de base
\[
\pm\delta\in
H^n_{\{x\}}\bigl(X_s,R\Psi_{\bar\eta}(\Lambda(m))\bigr).
\]

\pageoriginale

\subsubsection*{3.3.4.}
Volvamos a las hipótesis más generales de 3.3.1, con $n=2m+1$. Para
todo punto singular $x$ de $X_s$, denotaremos todavía por
$\pm\delta_x$ las secciones de $R^nf_!\Lambda(m)$ sobre $S$ que son
imágenes de los generadores del cociente primitivo de
$H^{n-1}(Y,\Lambda(m))$ (con la misma notación anterior) por el
morfismo compuesto
\begin{center}
\begin{tikzcd}[column sep=huge,row sep=huge]
H^{n-1}(Y,\Lambda(m))
  & H^{n-1}(\widetilde X_{s(x)},\Lambda(m))
      \arrow[l,"\sim"'] \arrow[r]
  & H^{n-1}(X^*_{s(x)},\Lambda(m)) \arrow[d] \\
H^0\bigl(S,R^nf_!\Lambda(m)\bigr)
  & H_c^n(X_s,\Lambda(m)) \arrow[l,"\sim"']
  & H^n_{\{x\}}(X_s,\Lambda(m)) \arrow[l]
\end{tikzcd}
\end{center}
La clase $\pm\delta_x$ es el \underline{ciclo evanescente} en $x$.

Cuando $S$ es un trait, la clase
\[
\pm\delta\in H_c^n(X_{\bar\eta},\Lambda(m))
\]
es la imagen, mediante XIII, 2.1.8.4, de la clase de 3.3.3
\[
\pm\delta\in
H^n_{\{x\}}\bigl(X_s,R\Psi_{\bar\eta}(\Lambda(m))\bigr).
\]

Demostraremos simultáneamente los dos resultados siguientes.

\subsubsection*{Proposición 3.3.5.}
\underline{Sea $S$ un trait henseliano y sea $X/S$ como en 3.3.3. La
variación viene dada por}
\[
\operatorname{Var}(\sigma)(a)
=(-1)^{m+1}\varepsilon_{b(x)}(\sigma)(a,\delta)\delta.
\]

\subsubsection*{Proposición 3.3.6.}
\underline{Sea $S$ estrictamente local y regular, y sea
$f:X\to S$ un morfismo propio y plano de dimensión relativa
$n=2m+1$ tal que $X_s$ tenga exactamente un punto no liso $x$, que se
supone cuadrático no degenerado. Supongamos además que $b(x)$ es un
parámetro. Entonces Galois actúa sobre la cohomología de la fibra
genérica geométrica mediante}
\[
\sigma(a)=a+(-1)^{m+1}
\varepsilon_{b(x)}(\sigma)(a,\delta_x)\delta_x.
\]

\medskip
\noindent\textup{(A)}
Sea $k$ un cuerpo algebraicamente cerrado de característica $0$, y sea
$S$ el henselizado

\pageoriginale

de $\operatorname{Spec}(k[T])$ en $0$. Demostremos 3.3.5 cuando $b(x)$
es un uniformizante.

Por el principio de Lefschetz, nos reducimos al caso $k=\mathbb C$.
Como la cuestión es local, podemos suponer que $f:X\to S$ se obtiene
por cambio de base de la aplicación
\[
\sum_i z_i^2:\mathbb C^{n+1}\longrightarrow\mathbb C.
\]
La construcción anterior de $\pm\delta$ tiene un análogo trascendente
en cohomología entera. Mediante XIV, 2.1, la clase $\pm\delta$
corresponde, por tanto, a un ciclo evanescente trascendente, imagen de
\[
\pm\delta_{\mathrm{cl}}
\in H^n_{\{0\}}\bigl(X_0,R\Psi_{\mathrm{cl}\eta}(\mathbb Z)\bigr).
\]
Por XIV, 2.1 y la propiedad correspondiente de los ciclos algebraicos
$\pm\delta$, la imagen de $\pm\delta_{\mathrm{cl}}$ en
\[
H^n_{\{0\}}\bigl(X_0,R\Psi_{\mathrm{cl}\eta}(\mathbb Z)\bigr)
\otimes\mathbb Z_\ell
\]
es un generador de este grupo. Así, $\pm\delta_{\mathrm{cl}}$ son los
dos generadores de
$H^n_{\{0\}}(X_0,R\Psi_{\mathrm{cl}\eta}(\mathbb Z))$, y coinciden con
los ciclos evanescentes de XIV, 3.2. Por consiguiente, la fórmula 3.3.5
resulta de XIV, 3.2.11. En rigor, habría que comprobar la compatibilidad
entre el isomorfismo de XIV, 2.1 y los productos cup y las trazas.

\medskip
\noindent\textup{(B)}
Demostremos 3.3.6 para $X/S$ como en (A), suponiendo además que $X/S$
es propio. Esto resulta de XIII, 2.4.6.2.

\medskip
\noindent\textup{(C)}
Sea $p$ un número primo, y sea $S_0$ el henselizado estricto de
$\operatorname{Spec}(\mathbb Z[T])$ en $(p,T)$. Demostremos 3.3.6 para
$X/S_0$ tal que $b(x)=T$.

La restricción del haz $R^if_*\Lambda$ al henselizado estricto de
$\operatorname{Spec}(\mathbb Z[T])$ en $(T)$ está regida por (B). Este
haz es localmente constante sobre el complemento del divisor liso
$T=0$, que tiene característica mixta. Por tanto, (C) resulta de (B) y
del lema de Abhyankar.

\par\medskip
\noindent\textup{(D)}
En al menos un caso al que se aplica (C), $\delta_x$ no es un elemento
de torsión en cohomología $\ell$-ádica.

\pageoriginale

Sea $S_1$ el henselizado estricto de $\operatorname{Spec}(\mathbb Z)$
en $(p)$, y sea $V\subset\mathbb P^r/S_1$ una variedad proyectiva lisa
sobre $S_1$, puramente de dimensión relativa $n+1$. Supongamos que
existe un pincel de Lefschetz de secciones hiperplanas para la fibra
especial $V_p$ de $V/S_1$ (XVII, 2.5). Este pincel está definido por
una recta $d'$ en la fibra especial del espacio proyectivo
$\check{\mathbb P}^{r}/S_1$ dual de
$\mathbb P^r/S_1$. Levantemos $d'$ a una recta $d$ de
$\check{\mathbb P}^{r}/S_1$, y sea $X\to d$ el pincel de secciones
hiperplanas de $V$ parametrizadas por $d$ (XVIII, 3.1). El lugar
excepcional $\Delta$ de $d$ es étale sobre $S_1$ (XVII, 6). Tras una
localización étale en un punto de $\Delta_p\subset d$, se obtiene, por
tanto, una situación de tipo (C).

Para comprobar que $\delta_x$ no es de torsión, basta hacerlo tras
extender los escalares de $S_1$ a $\mathbb C$. Tomemos
$V=\mathbb P^{n+1}$ con la inmersión proyectiva definida por
$\mathcal O(N)$. Para $N\ge2$, se aplica XVII, 2.5 (en este caso
particular es válido para todo $N$). Utilizamos entonces el lema
siguiente.

\subsubsection*{Lema 3.3.7.}
\underline{Para un pincel de Lefschetz de hipersuperficies de grado
$N$ suficientemente grande en $\mathbb P^{n+1}(\mathbb C)$ (basta
$N\ge3$), un ciclo evanescente en cohomología $\mathbb Q_\ell$ no es
nulo.}

Este lema resulta de la teoría de los pinceles de Lefschetz, por ejemplo
de XVIII, 6.6.1. No hay circularidad al utilizar aquí este resultado de
XVIII, basado en la ``fórmula de Picard--Lefschetz'', puesto que ya se
dispone de (A). También puede utilizarse la variante más directa y
puramente trascendente XIX, n\textsuperscript{o}~4, que solo depende de
XIV, 3.2.

\medskip
\noindent\textup{(E)}
Demostremos 3.3.5 en característica residual $p$ para
$(X/S,x\in X_s)$.

Sea $X_0/S_0$ como en (D). Por 1.3.2(i), existe un morfismo
$g:S\to S_0$ tal que el henselizado en $x_0$ de
\[
g^*X_0=X_0\times_{S_0}S
\]
es $S$-isomorfo al henselizado de $X$ en $x$. Basta, por tanto,
demostrar 3.3.5 para $g^*X_0$. Por cambio de base a partir de (C) para
$X_0/S_0$, se obtiene
\[
\sigma(a)=a-(-1)^m
\varepsilon_{b(x)}(\sigma)(a,\delta_x)\delta_x
\]

\pageoriginale

en $H^n(X_{0\bar\eta})$, para $\sigma$ en el grupo de inercia. Sea
$\pi$ un uniformizante y comparemos esta fórmula con (2.2.5.9), donde
$\varepsilon=\varepsilon_\pi$. Se obtiene
\[
\lambda_X\varepsilon_\pi(\sigma)(a,\delta_x)\delta_x
=-(-1)^m\varepsilon_{b(x)}(\sigma)(a,\delta_x)\delta_x
\]
en $H^n(X_{0\bar\eta})$. Como $\delta_x$ no es de torsión, podemos
cancelar $(a,\delta_x)\delta_x$ en esta identidad. La fórmula
(2.2.5.9), con el valor resultante de $\lambda_X$, es precisamente
3.3.5.

\medskip
\noindent\textup{(F) Demostración de 3.3.6 en el caso general.}

El caso general de 3.3.6 se reduce al caso general de 3.3.5 del mismo
modo que (B) y (C) se reducen a (A).

\subsection*{3.4. Resumen}

Las Proposiciones 3.1.2, 3.2.1, 3.2.2, 3.2.3 y 3.3.5 constituyen la
``fórmula de Picard--Lefschetz''. Esta fórmula es el análogo algebraico
de XIV, 3.2.11, pero adopta formas muy distintas según la paridad de la
dimensión relativa.

En el caso más interesante, en que $X/S$ es propio, los resultados se
resumen como sigue.

\subsubsection*{Teorema 3.4 (fórmula de Picard--Lefschetz).}
\underline{Sea $S$ un trait estrictamente henseliano y sea
$f:X\to S$ un morfismo propio y plano, puramente de dimensión relativa
$n=2m$ o $2m+1$. Supongamos que $X_\eta$ es liso y que $X_s$ solo tiene
singularidades cuadráticas ordinarias. Sea $\Sigma$ el lugar singular
de $X_s$.}

\begin{enumerate}
\item[(i)] \underline{A cada $x\in\Sigma$ se asocia un ciclo
evanescente}
\[
\pm\delta_x\in H_n(X_{\bar\eta})(m).
\]
\underline{Los $\delta_x$ son ortogonales dos a dos. Cuando $n$ es par,}
\[
(\delta_x,\delta_x)=(-1)^m\,2.
\]

\item[(ii)] \underline{Existen isomorfismos}
\[
H^i(X_s)\xrightarrow{\ \sim\ }H^i(X_{\bar\eta})
\]
\underline{para $i\ne n,n+1$, y una sucesión exacta}

\pageoriginale

{\small
\[
0\longrightarrow H^n(X_s)
 \longrightarrow H^n(X_{\bar\eta})
 \xrightarrow{\,((\,\cdot\,,\delta_x))_{x\in\Sigma}\,}
 \bigoplus_{x\in\Sigma}\Lambda(n-m)
 \longrightarrow H^{n+1}(X_s)
 \longrightarrow H^{n+1}(X_{\bar\eta})
 \longrightarrow0.
\]
}

\item[(iii)] \underline{Se definen caracteres locales del grupo de
inercia $I$:}
\[
\begin{array}{lll}
n\text{ par}:&\varepsilon_x:I\longrightarrow\{\pm1\},\\[2pt]
n\text{ impar}:&\varepsilon_x:I\longrightarrow\Lambda(1).
\end{array}
\]
\underline{Cuando $X$ es regular en $x$, estos caracteres son
sobreyectivos; cuando $x$ es no degenerado, son, respectivamente, el
único carácter de orden $2$ y el carácter canónico.}

\underline{La acción del grupo de inercia $I$ viene dada por}
\[
n=2m\qquad
\sigma(a)=a+(-1)^m\sum_{x\in\Sigma}
 \frac{\varepsilon_x(\sigma)-1}{2}
 \sisi{(a\delta_x)}{(a,\delta_x)}\delta_x,
\]
\[
n=2m+1\qquad
\sigma(a)=a-(-1)^m\sum_{x\in\Sigma}
 \varepsilon_x(\sigma)\sisi{(a\delta_x)}{(a,\delta_x)}\delta_x.
\]
\underline{En particular,}
\[
\begin{array}{lll}
n\text{ par}:&\sigma(\delta_x)=\varepsilon_x(\sigma)\delta_x,\\[2pt]
n\text{ impar}:&\sigma(\delta_x)=\delta_x.
\end{array}
\]
\end{enumerate}
